% !TEX root = ../../main_power.tex
% ---------------------------------------------------------------------
%  Kinetis K22 / tinyK22 (Teil 5)
% ---------------------------------------------------------------------
\section{Kinetis und tinyK22}
\label{sec:kinetis}

\begin{falle}[Sicherheitsgurt -- zuerst lesen]
Geht der K22 sofort nach dem Start in Reset oder Fault, ist der
K20-Debug-Chip (Boards vor V1.3) nicht schnell genug, um das Target
zwischen zwei Resets anzuhalten \ra \textbf{kein Zugriff mehr auf den
K22}. Jedes Low-Power-Experiment beginnt deshalb mit einer Verzögerung,
die dem Debugger eine Chance gibt:
\begin{lstlisting}
static void safetyBelt(void) {
  for(int i=0; i<20000000; i++) {
    __asm("nop"); /* give debugger a chance */
  }
}
int main(void) { safetyBelt(); ... }
\end{lstlisting}
\end{falle}

\subsection{Board}
\begin{ztab}{@{}lX@{}}
Speisung & 5\,V über K20-USB oder K22-USB \ra \R{VREGIN} \ra
  \R{VOUT} 3.3\,V, insgesamt max.\ 120\,mA. Schottky-Diode
  \R{PMEG2005CT} gegen Rückstrom; \R{JP1} setzen, wenn der K20 nicht
  gebraucht wird \\
SD-Karte & optional, mit eigenem Regler \R{TC1262-3.3VDB} (500\,mA,
  Wirkungsgrad 80\,\%). Low Power: industrielle Karten, wenig lesen und
  vor allem \emph{wenig schreiben} \\
Pins & Debug-/JTAG-Pins haben in der Default-Konfiguration teils
  Pull-ups. Sonst: Default-Muxing lassen und floating -- siehe
  „Configuration for Unused Pins/GPIOs on Kinetis“ \\
\end{ztab}

\subsection{Flash zur Laufzeit programmieren}
Gebraucht für Bootloader, Kalibrierdaten, Einstellungen und Logs.
Programmieren erfordert \T{höhere interne Spannungen} aus Ladungspumpen
\ra deutlich höherer Verbrauch und Anforderungen an die Speisung.
\begin{falle}[Nicht in jedem Modus]
Flash-Programmierung ist \emph{nicht} in allen Power-Modi und bei jeder
Taktfrequenz erlaubt -- aus \R{HSRUN} muss man vorher zurück in
\R{RUN}, und Interrupts werden abgeschaltet:
\begin{lstlisting}
SMC_SetPowerModeRun(SMC);   /* Fcore=60MHz, Fbus=Fflash=20MHz */
McuWait_Waitms(1);          /* Takt umschalten lassen */
uint32_t primask = DisableGlobalIRQ();
FLASH_Erase(&s_flashDriver, FLASH_NVM_ADDR_START, ...);
\end{lstlisting}
\end{falle}

\subsection{Clock Gating und Gear Shifting}
\label{rez:gearshift}
\begin{lstlisting}
CLOCK_EnableClock(kCLOCK_PortA);   /* Clock enabled  */
CLOCK_DisableClock(kCLOCK_PortA);  /* Clock disabled */
\end{lstlisting}
\begin{falle}[HardFault ohne Takt]
Ein Zugriff auf eine Peripherieeinheit, die \emph{nicht} mit einem Takt
versorgt ist, löst einen \T{HardFault} aus. Beim Suchen eines
HardFaults ist ein vergessenes \R{CLOCK\_EnableClock} ein
Standardverdächtiger.
\end{falle}
\SW{Gear Shifting} Die Clock-\emph{Modi} zu wechseln ist kompliziert und
dauert lange -- in der Praxis ändert man die \T{Clock-Divider}:
\begin{lstlisting}
configFast.clkdiv1 = 0x44440000U;  /* 20 MHz */
configSlow.clkdiv1 = 0x99990000U;  /* 10 MHz */
CLOCK_SetSimConfig(slow ? &configSlow : &configFast);
SystemCoreClock = CLOCK_GetCoreSysClkFreq();
UpdateUartClock();                 /* nicht vergessen! */
\end{lstlisting}
Die aktuelle Frequenz steht in \R{SystemCoreClock}. Beim Umschalten muss
sich der Takt \emph{stabilisieren} (Zeit einrechnen), und alles, was den
Takt mitbenutzt, muss nachgezogen werden -- der UART sonst sendet
Buchstabensalat.

\subsection{ARM Power Modes}
\label{rez:modi}
\begin{formel}[Drei Prinzipien, viele Modi]
Jeder Modus kombiniert \T{Clock Gating}, abgeschaltete
\T{Power Domains} und eine zulässige \T{Wakeup-Quelle}. Die Quelle
unterscheidet die Klassen: \emph{synchron} braucht noch einen Takt,
\emph{asynchron} kommt ohne aus.
\end{formel}
\begin{ztab}{@{}llX@{}}
\T{Run} & \R{HSRUN RUN VLPR} & Core läuft mit dem Core-Clock, alles
  aktiv \\
\T{Sleep} & \R{WAIT VLPW} & Core-Clock gated (\R{wfi}), Peripherie
  aktiv, der \T{NVIC} weckt \emph{synchron} \\
\T{Deepsleep} & \R{STOP VLPS} \R{LLS2 LLS3} \R{VLLS3 VLLS2} \R{VLLS1
  VLLS0} & Aufwecken nur \emph{asynchron} über \T{AWIC} (Asynchronous
  Wakeup Interrupt Controller) oder \T{LLWU} (Low Leakage Wakeup
  Unit) \\
\end{ztab}
\merkzeile \R{VLP\dots} heisst Very Low Power, \R{LLS} Low Leakage Stop,
\R{VLLS} Very Low Leakage Stop -- je weiter unten in der Liste, desto
weniger bleibt erhalten und desto länger dauert das Aufwachen.
